% !TeX spellcheck = ru_RU
% !TEX root = vkr.tex

\section{Особенности реализации}

\subsection{Используемый стек технологий}
\subsection{Симуляция}
\subsection{Реализованные алгоритмы}

В рамках работы в систему интегрированы следующие мультиагентные алгоритмы определения позиций объектов.

\subsubsection{Алгоритм на основе SPSA}


\[
\left\{
\begin{aligned}
  &\mathbf{x}_{2k}^i = \hat{\boldsymbol{\theta}}_{2k-2}^i + \beta_k^+ \boldsymbol{\Delta}_k^i, \\[0.5em]
  &\mathbf{x}_{2k-1}^i = \hat{\boldsymbol{\theta}}_{2k-2}^i - \beta_k^- \boldsymbol{\Delta}_k^i \\[0.5em]
  &\hat{\boldsymbol{\theta}}_{2k-1}^i = \hat{\boldsymbol{\theta}}_{2k-2}^i, \\[0.5em]
  &\hat{\boldsymbol{\theta}}_{2k}^i = \hat{\boldsymbol{\theta}}_{2k-1}^i - \alpha \left( \frac{y_{2k}^i - y_{2k-1}^i}{\beta_k} \mathbf{K}_k^i(\boldsymbol{\Delta}_k^i) \right. \\[0.5em]
  &\qquad \left. {} + \gamma \sum_{j \in \mathcal{N}_{2k-1}^i} b_{2k-1}^{i,j} (\tilde{\theta}_{2k-1}^{i,j} - \hat{\theta}_{2k-1}^i) \right)
\end{aligned}
\right.
\]


\paragraph{Пояснение основных шагов.}

\begin{enumerate}
    \item \textbf{Формирование текущей оценки.}  
    Агент использует предыдущее состояние $\hat{\theta}_{2k-2}^i$ как базовую точку $x_k^i$.

    \item \textbf{SPSA-возмущение.}  
    Генерируется случайный вектор $\Delta_k^i \in \{\pm 1\}^d$ и формируются две точки:
    $x_{2k}^i$ и $x_{2k-1}^i$.

    \item \textbf{Оценка градиента.}  
    Производится симметричная разностная аппроксимация:
    \[
    \frac{y_{2k}^i - y_{2k-1}^i}{\beta_k}\Delta_k^i.
    \]

    \item \textbf{Консенсусный шаг.}  
    Добавляется вклад соседей через граф взаимодействия:
    \[
    \sum_{j \in \mathcal{N}_i} b^{i,j} (\hat{\theta}^j - \hat{\theta}^i).
    \]

    \item \textbf{Обновление параметров.}  
    Выполняется градиентный шаг:
    \[
    \hat{\theta}_{2k}^i = \hat{\theta}_{2k-1}^i - \alpha g_k^i.
    \]
\end{enumerate}

\subsubsection{Ускоренный алгоритм на основе SPSA + Консенсус}


\[
\left\{
\begin{aligned}
  &\bar{\mathbf{x}}_k = \dfrac{1}{\gamma_{k-1} + \alpha_k(\mu - \eta)} \Big( \alpha_k \gamma_{k-1} \bar{\mathbf{z}}_{k-1} + \gamma_k \bar{\boldsymbol{\theta}}_{2k-1} \Big), \\[0.8em]
  &\bar{\mathbf{x}}_{2k} = \bar{\mathbf{x}}_k + \beta_k^+ \bar{\boldsymbol{\Delta}}_k, \quad
  \bar{\mathbf{x}}_{2k-1} = \bar{\mathbf{x}}_k - \beta_k^- \bar{\boldsymbol{\Delta}}_k, \\[0.8em]
  &\bar{\mathbf{g}}_k = \left( \dfrac{\bar{\mathbf{y}}_{2k} - \bar{\mathbf{y}}_{2k-1}}{\beta_k} \otimes \mathbf{I}_d \right) \bar{\boldsymbol{\Delta}}_k + \omega \left( \bar{\mathcal{L}}_{2k-1} \otimes E_n \right) \bar{\mathbf{x}}_k, \\[0.8em]
  &\bar{\boldsymbol{\theta}}_{2k} = \bar{\mathbf{x}}_k - h \bar{\mathbf{g}}_k, \\[0.5em]
  &\bar{\boldsymbol{\theta}}_{2k+1} = \bar{\boldsymbol{\theta}}_{2k}, \\[0.5em]
  &\bar{\mathbf{z}}_k = \dfrac{(1 - \alpha_k) \gamma_{k-1} \bar{\mathbf{z}}_{k-1} + \alpha_k (\mu - \eta) \bar{\mathbf{x}}_k - \alpha_k \bar{\mathbf{g}}_k}{\gamma_k}
\end{aligned}
\right.
\]

\paragraph{Пояснение основных шагов.}

\begin{enumerate}
    \item \textbf{Формирование точки зондирования $\tilde{x}_k$.}  
    Это взвешенная комбинация предыдущего состояния $\theta$ и вспомогательной переменной $z_k$.  
    Она играет роль ``ускоренной'' точки, аналогично методам Нестерова.

    \item \textbf{Случайное возмущение (SPSA).}  
    Генерируется случайный вектор $\Delta_k$ (обычно Бернулли $\pm 1$).  
    Затем вычисляются две точки:
    \[
    x_{2k},\quad x_{2k-1}
    \]
    что позволяет оценить градиент без явных производных.

    \item \textbf{Оценка градиента.}  
    Градиент аппроксимируется как
    \[
    g_k \approx \frac{f(x+\beta\Delta) - f(x-\beta\Delta)}{2\beta}\Delta,
    \]
    с добавлением консенсусного члена $\omega L x_k$, учитывающего взаимодействие агентов.

    \item \textbf{Шаг спуска.}  
    Выполняется обновление:
    \[
    \theta_{2k} = \tilde{x}_k - h g_k,
    \]
    где $h$ — шаг обучения.

    \item \textbf{Ускорение (переменная $z_k$).}  
    Вводится вспомогательная последовательность $z_k$, которая аккумулирует информацию о предыдущих итерациях и обеспечивает ускоренную сходимость.

    \item \textbf{Обновление параметров ускорения.}  
    Параметры $\alpha_k$ и $\gamma_k$ обновляются рекурсивно:
    \[
    \gamma_k = (1-\alpha_k)\gamma_{k-1} + \alpha_k(\mu - \eta),
    \]
    что задаёт динамику ускорения.
\end{enumerate}